\section{Value Learning：Deep Q-Networks (DQN)}

\subsection{DQN 的基本思想}

Deep Q-Network (DQN) 是 Value Learning 的代表性算法。其核心思想是：用一个深度神经网络来\textbf{近似 Q-function}，然后利用学到的 Q-function 来推导最优策略。

\begin{figure}[H]
\centering
\includegraphics[width=0.95\textwidth]{slides-images/slide-26.jpg}
\caption{DQN 的两种架构：左侧为 (state, action) $\to$ Q-value，右侧为 state $\to$ 所有动作的 Q-value\protect\footnotemark}
\end{figure}
\footnotetext{Slides 第26页。}

DQN 有两种等价的架构设计：

\begin{enumerate}
    \item \textbf{方案一}：输入 $(s, a)$，输出单个 $Q(s,a)$ 值。需要对每个可能的动作调用一次网络。
    \item \textbf{方案二（更高效）}：输入 $s$，同时输出所有动作的 Q 值 $Q(s, a_1), Q(s, a_2), \ldots, Q(s, a_n)$。只需一次前向传播。
\end{enumerate}

在实践中通常采用方案二，因为它更高效------一次网络调用即可获得所有动作的价值评估。

\subsection{DQN 的决策过程}

\begin{figure}[H]
\centering
\includegraphics[width=0.85\textwidth]{slides-images/slide-36.jpg}
\caption{DQN 决策示例：选择 Q 值最高的动作 $a_1$（向左移动）\protect\footnotemark}
\end{figure}
\footnotetext{Slides 第36页。}

决策过程非常直接：
\begin{enumerate}
    \item 将当前状态 $s$（如 Atari 游戏画面）输入 DQN
    \item 网络输出每个可能动作的 Q 值
    \item 选择 Q 值最大的动作执行：$\pi(s) = \arg\max_a Q(s, a)$
\end{enumerate}

\subsection{Atari Breakout 示例：理解 Q 值}

课上通过 Atari Breakout 游戏来深入理解 Q-function 的含义。考虑两种 state-action pair：

\begin{figure}[H]
\centering
\includegraphics[width=0.95\textwidth]{slides-images/slide-22.jpg}
\caption{Atari Breakout 中的 Q 值比较：哪个 (s, a) 对的 Q 值更高？\protect\footnotemark}
\end{figure}
\footnotetext{Slides 第22页。}

\begin{itemize}
    \item \textbf{场景 A}：球直直落向球拍正中央，球拍保持不动
    \item \textbf{场景 B}：球飞向侧面，球拍向右移动追赶
\end{itemize}

\begin{knowledgebox}{Q 值的直觉理解}
直觉上场景 B 的 Q 值可能更高（球弹向侧面打到更多砖块），但实际实验表明，场景 B 的策略（侧击）虽然单次得分可能更高，但球拍移动到侧边后回球速度变慢，长期来看反而不利。这说明 Q 值衡量的是\textbf{长期累积回报}而非单步奖励，人类直觉可能无法准确估计 Q 值------这正是需要深度网络来学习的原因。
\end{knowledgebox}

\subsection{DQN 的训练}

DQN 的训练核心是定义合适的目标值和损失函数。

\textbf{Target}（目标值）：利用 Bellman 方程构造训练目标：

$$\text{target} = r + \gamma \max_{a'} Q(s', a')$$

\begin{itemize}
    \item $r$：当前步获得的即时奖励
    \item $s'$：执行动作后到达的新状态
    \item $\max_{a'} Q(s', a')$：在新状态下，所有可能动作中最大的 Q 值（假设后续采取最优策略）
\end{itemize}

\textbf{Q-Loss}（损失函数）：

$$\mathcal{L} = \mathbb{E}\left[\left\| \underbrace{\left(r + \gamma \max_{a'} Q(s', a')\right)}_{\text{target}} - \underbrace{Q(s, a)}_{\text{predicted}} \right\|^2\right]$$

\begin{figure}[H]
\centering
\includegraphics[width=0.95\textwidth]{slides-images/slide-30.jpg}
\caption{DQN 的训练：Q-Loss 定义为 target 与 predicted Q 值之间的均方误差\protect\footnotemark}
\end{figure}
\footnotetext{Slides 第30页。}

\begin{importantbox}{DQN 训练的核心直觉}
DQN 的训练本质上是让网络预测的 Q 值逐渐逼近由 Bellman 方程给出的 target 值。target 由"即时奖励 + 折扣后的最优未来回报"构成，这就实现了"用未来回报来指导当前决策"的目标。训练过程通过最小化 predicted Q 与 target Q 之间的均方误差来进行反向传播和参数更新。
\end{importantbox}

\subsection{DQN 在 Atari 游戏上的成果}

Google DeepMind 于 2015 年在 Nature 上发表了 DQN 的里程碑论文，使用相同的网络架构和超参数在 49 种 Atari 游戏上进行测试。

\begin{figure}[H]
\centering
\includegraphics[width=0.95\textwidth]{slides-images/slide-32.jpg}
\caption{DQN 的网络架构：卷积层处理游戏画面，全连接层输出每个动作的 Q 值\protect\footnotemark}
\end{figure}
\footnotetext{Slides 第32页。Mnih+ Nature 2015.}

网络架构包含：
\begin{itemize}
    \item 多层 \textbf{Convolution}（卷积层）：处理原始游戏画面像素
    \item 多层 \textbf{Fully Connected}（全连接层）：将特征映射到各个动作的 Q 值
    \item 输出层：每个节点对应一个可能的动作（如游戏手柄的各个方向和按键组合）
\end{itemize}

\begin{figure}[H]
\centering
\includegraphics[width=0.95\textwidth]{slides-images/slide-33.jpg}
\caption{DQN 在 Atari 游戏上的表现：超过 50\% 的游戏达到或超越人类水平\protect\footnotemark}
\end{figure}
\footnotetext{Slides 第33页。Mnih+ Nature 2015.}

\begin{knowledgebox}{DQN 的历史意义}
DQN 的重要性在于它证明了一个\textbf{端到端}的深度强化学习系统可以直接从原始像素学会复杂的策略。在此之前，RL 通常需要手工设计特征（Feature Engineering）。DQN 开启了 Deep RL 的新时代，也是 DeepMind 被 Google 收购的关键技术之一。
\end{knowledgebox}

\subsection{Q-Learning 的局限性}

\begin{figure}[H]
\centering
\includegraphics[width=0.85\textwidth]{slides-images/slide-34.jpg}
\caption{Q-Learning 的两大局限：无法处理连续动作空间，且只能学习确定性策略\protect\footnotemark}
\end{figure}
\footnotetext{Slides 第34页。}

尽管 DQN 取得了巨大成功，Q-Learning 方法存在两个重要局限：

\textbf{复杂度问题}（Complexity）：
\begin{itemize}
    \item 只能处理\textbf{离散且有限}的动作空间
    \item 对于连续动作空间（如控制方向盘的转角），无法枚举所有可能动作
\end{itemize}

\textbf{灵活性问题}（Flexibility）：
\begin{itemize}
    \item 策略是通过 $\arg\max$ 确定性地从 Q-function 推导的
    \item 无法学习\textbf{随机策略}（Stochastic Policy），而某些场景下随机性是有益的
\end{itemize}

\begin{warningbox}{何时 Q-Learning 不适用？}
当动作空间是连续的（如机器人关节角度、汽车油门刹车力度），Q-Learning 无法直接应用，因为它需要对\textbf{所有可能动作}计算 Q 值并取最大值。在连续空间中动作有无穷多个，这在计算上不可行。此时需要转向 Policy Gradient 方法。
\end{warningbox}

\subsection{本章小结}

DQN 通过深度神经网络近似 Q-function，实现了从原始像素到决策的端到端学习。其训练基于 Bellman 方程构造 target，通过最小化 Q-Loss 来更新网络参数。DQN 在 Atari 游戏上取得了超越人类的成绩，但受限于离散动作空间和确定性策略，无法直接处理连续控制等更复杂的问题。


